% Options for packages loaded elsewhere
\PassOptionsToPackage{unicode}{hyperref}
\PassOptionsToPackage{hyphens}{url}
\PassOptionsToPackage{dvipsnames,svgnames,x11names}{xcolor}
%
\documentclass[
  11pt,
  a4paper,
]{ltjsarticle}
\usepackage{amsmath,amssymb}
\usepackage{iftex}
\ifPDFTeX
  \usepackage[T1]{fontenc}
  \usepackage[utf8]{inputenc}
  \usepackage{textcomp} % provide euro and other symbols
\else % if luatex or xetex
  \usepackage{unicode-math} % this also loads fontspec
  \defaultfontfeatures{Scale=MatchLowercase}
  \defaultfontfeatures[\rmfamily]{Ligatures=TeX,Scale=1}
\fi
\usepackage{lmodern}
\ifPDFTeX\else
  % xetex/luatex font selection
\fi
% Use upquote if available, for straight quotes in verbatim environments
\IfFileExists{upquote.sty}{\usepackage{upquote}}{}
\IfFileExists{microtype.sty}{% use microtype if available
  \usepackage[]{microtype}
  \UseMicrotypeSet[protrusion]{basicmath} % disable protrusion for tt fonts
}{}
\makeatletter
\@ifundefined{KOMAClassName}{% if non-KOMA class
  \IfFileExists{parskip.sty}{%
    \usepackage{parskip}
  }{% else
    \setlength{\parindent}{0pt}
    \setlength{\parskip}{6pt plus 2pt minus 1pt}}
}{% if KOMA class
  \KOMAoptions{parskip=half}}
\makeatother
\usepackage{xcolor}
\usepackage[margin=24mm]{geometry}
\setlength{\emergencystretch}{3em} % prevent overfull lines
\providecommand{\tightlist}{%
  \setlength{\itemsep}{0pt}\setlength{\parskip}{0pt}}
\setcounter{secnumdepth}{5}
\ifLuaTeX
  \usepackage{selnolig}  % disable illegal ligatures
\fi
\IfFileExists{bookmark.sty}{\usepackage{bookmark}}{\usepackage{hyperref}}
\IfFileExists{xurl.sty}{\usepackage{xurl}}{} % add URL line breaks if available
\urlstyle{same}
\hypersetup{
  colorlinks=true,
  linkcolor={blue},
  filecolor={Maroon},
  citecolor={Blue},
  urlcolor={blue},
  pdfcreator={LaTeX via pandoc}}

\author{}
\date{}

\begin{document}

\hypertarget{ux5b66ux7fd2ux7d4cux8def}{%
\section{学習経路}\label{ux5b66ux7fd2ux7d4cux8def}}

\hypertarget{stage-1-ux7a7aux9593ux3092ux7406ux89e3ux3059ux308b}{%
\subsection{Stage 1:
空間を理解する}\label{stage-1-ux7a7aux9593ux3092ux7406ux89e3ux3059ux308b}}

\texttt{01\_vectors\_and\_spaces}
では、ベクトルを矢印としてだけでなく、加法とスカラー倍ができる対象として捉えます。

到達目標は、\texttt{span}、一次独立、基底、次元を自分の言葉で説明できることです。

\hypertarget{stage-2-ux884cux5217ux3092ux5199ux50cfux3068ux3057ux3066ux7406ux89e3ux3059ux308b}{%
\subsection{Stage 2:
行列を写像として理解する}\label{stage-2-ux884cux5217ux3092ux5199ux50cfux3068ux3057ux3066ux7406ux89e3ux3059ux308b}}

\texttt{02\_matrices\_and\_linear\_maps}
では、行列積を計算規則ではなく写像の合成として捉えます。

特に

\[
\dim\ker A+\operatorname{rank}(A)=n
\]

が「入力の自由度が、消える方向と残る方向に分かれる」という意味を持つことを理解します。

\hypertarget{stage-3-ux65b9ux7a0bux5f0fux3092ux7a7aux9593ux3068ux3057ux3066ux7406ux89e3ux3059ux308b}{%
\subsection{Stage 3:
方程式を空間として理解する}\label{stage-3-ux65b9ux7a0bux5f0fux3092ux7a7aux9593ux3068ux3057ux3066ux7406ux89e3ux3059ux308b}}

\texttt{Ax=b} は数値を求める問題であると同時に、\texttt{b}
が列空間に入っているかを問う問題です。

Gauss
消去を「アルゴリズム」としてだけでなく、「解空間の構造を露出させる操作」として学びます。

\hypertarget{stage-4-ux5e7eux4f55ux5b66ux3092ux5165ux308cux308b}{%
\subsection{Stage 4:
幾何学を入れる}\label{stage-4-ux5e7eux4f55ux5b66ux3092ux5165ux308cux308b}}

内積を導入すると、長さ・角度・直交・最短距離を定義できます。ここから射影と最小二乗が自然に出てきます。

\hypertarget{stage-5-ux7ddaux5f62ux5909ux63dbux306eux4e0dux5909ux306aux65b9ux5411ux3092ux63a2ux3059}{%
\subsection{Stage 5:
線形変換の不変な方向を探す}\label{stage-5-ux7ddaux5f62ux5909ux63dbux306eux4e0dux5909ux306aux65b9ux5411ux3092ux63a2ux3059}}

固有値・固有ベクトルは

\[
Av=\lambda v
\]

を満たす特別な方向です。ここで「良い基底を選ぶと変換が単純になる」という考え方を身につけます。

\hypertarget{stage-6-ux5206ux89e3ux3059ux308b}{%
\subsection{Stage 6: 分解する}\label{stage-6-ux5206ux89e3ux3059ux308b}}

LU、QR、固有値分解、SVD を比較します。特に SVD
は一般の長方形行列にも使え、ランク・最小二乗・擬似逆・低ランク近似を統一します。

\hypertarget{ux5b66ux7fd2ux65b9ux6cd5}{%
\subsection{学習方法}\label{ux5b66ux7fd2ux65b9ux6cd5}}

各節で次の順序を守ることを推奨します。

\begin{enumerate}
\def\labelenumi{\arabic{enumi}.}
\tightlist
\item
  定義を読む。
\item
  図形的・構造的な意味を言葉で説明する。
\item
  数式を自力で導出する。
\item
  小さな行列で手計算する。
\item
  演習を解く。
\item
  NumPy で検算する。
\end{enumerate}

計算だけ解けても、定義の意味を説明できなければ次に進まないのが理想です。

\end{document}
